\begin{conf-abstract}
{Establishment of a new Groundwater Dating Laboratory at TU Dresden}
{Diana Burghardt, Paul Fremdling}
{TU Dresden, Institut für Grundwasserwirtschaft}
{Groundwater is the most important source of drinking and industrial water worldwide. Groundwater models and average residence times of groundwater are needed to assess usable groundwater reserves and their sustainable management. The only reliable method for dating sustainably usable groundwater is the analysis of helium isotopes dissolved in the water. Funding from the EFRE-InfraProNet 2021-2027 support programme and the Free State of Saxony will be used to establish a laboratory for groundwater dating using helium isotope analysis at the Institute for Groundwater Management at TU Dresden from 2026 to 2027. In addition to the funds for the system of all devices, personnel funds are also available. The objectives of the funding are (i) the technological establishment of this laboratory, (ii) the verification of helium isotope analysis of the new system in comparison to that of the helis laboratory at the University of Bremen, (iii) the validation of helium isotope analysis of the new system on representative groundwater samples in Saxony, and (iv) the development of a web map of sustainably usable groundwater reserves in Saxony. }
\end{conf-abstract}
